\section{来源审计与问题边界：这是一堂 2023 年的对齐课}

本讲录制于 2023 年 1 月 17 日，处在 ChatGPT 刚公开不久的时间点。它讨论的是当时已经形成的 RLHF 路线、正在探索的 scalable oversight，以及学生围绕偏好、欺骗、工具、inner alignment 与 interpretability 提出的开放问题。旧稿把后来出现的 ``Superalignment'' 组织叙事、通用上线流程和产品治理清单倒灌进课堂，因此本次按官方视频、1,264 条人工字幕、17 张教学画面和六篇当时已公开的一手论文完整重写，而不是在旧结构上补丁式扩写。

\lecturefigure{slide-01-title.jpg}{Jan Leike 的讲座标题：Aligning Language Models with Humans。}{Stanford Online 官方视频 00:00:23--00:00:49。}

标题中的 ``with humans'' 比 ``make models safe'' 更具体：本讲要问的是，人类怎样把无法完整写成规则的意图转化为训练信号，并在模型能力持续提升时保持这种信号可用。读者应把三类证据分开：slide 给出公开主线，讲者口头解释补充动机与限制，一手论文提供公式和实验细节；三者不一致时，讲义会明确标注边界，而不是把后来的知识假装成课堂原话。

\begin{importantbox}{本讲的证据合同}
课堂中的数字、系统状态和研究判断都以 2023-01-17 为界。InstructGPT 的公开论文可以用于澄清 SFT、reward model 与 PPO 的关系；2023 年之后的项目名称、部署架构和政策变化不能被写成讲者当时已经给出的答案。后半小时几乎一直显示最后一张 slide，因此 teacher voice 主要来自时间戳字幕，而不是重复插入同一人物画面。
\end{importantbox}

\subsection{Team AI：为什么能力增长会把对齐变成战略问题}

讲者先不用损失函数，而是构造一场 ``Team AI 对 Team Human'' 的比赛。这个类比的重点不是断言当前系统已经全面超过人类，而是把动态性带入问题：不同 AI player 会逐步入场，后来的 player 可能更快、更便宜、在更多任务上更强；人类眼下仍有一个重要优势---决定哪些系统以什么方式进入真实世界，并训练其中一些系统站到人类这一边。

\lecturefigure{slide-02-team-ai-strong-players.jpg}{Team AI 正在加入比赛，而且未来可能出现极强的 player。}{Stanford Online 官方视频 00:00:49--00:03:18。}

\begin{knowledgebox}{读图：强 player 不是“今天已经无所不能”}
这页把模型能力看成随时间加入的 player 序列。讲者同时强调 2023 年系统仍有明显局限；ChatGPT 只是一个锚点，展示知识广度、语言覆盖和低边际成本可能怎样改变分工。图中 ``incredibly strong'' 是面向未来的风险建模假设，不是对任意当代模型的实测结论。
\end{knowledgebox}

这个类比随后分成两个行动目标。第一，训练或选择愿意帮助人类、按人类意图行动的 AI player；第二，设计制度和规则，使拥有强 AI 的竞争环境不把人类推向失败。第一项被讲者称为 alignment，第二项更接近 governance。二者互相影响，却不是同一个技术问题。

\lecturefigure{slide-03-two-objectives.jpg}{Team Human 的两个目标：招募 AI player，以及制定不会输掉比赛的规则。}{Stanford Online 官方视频 00:03:18--00:04:16。}

\begin{warningbox}{不要把 alignment 与 governance 混成一个万能词}
Alignment 研究训练信号、行为目标、评估与泛化；governance 处理部署权限、竞争规则、责任与制度协调。本讲只展开前者。把组织级流程、监管方案或市场竞争直接塞进后文，不但偏离课堂，也会掩盖本讲真正的技术瓶颈：当人类看不懂模型完成的任务时，怎样继续给出可信反馈。
\end{warningbox}

\subsection{本章小结}

本章建立了两条边界。时间上，它是一堂 2023 年初的课；问题上，它聚焦如何让模型追随 human intent，而不试图解决所有 AI governance。Team AI 类比提供战略动机，接下来要把 ``站在人类一边'' 从口号拆成可训练目标。

